% SECCIÓN 5: EL DAÑO MÉDICO — TIPOLOGÍA Y CRITERIOS JURISPRUDENCIALES DE VALORACIÓN
%
% Objetivos de esta sección:
%   1. Presentar la tipología del daño en el ámbito médico:
%      - Daño patrimonial: daño emergente (gastos médicos, rehabilitación) y
%        lucro cesante (pérdida de ingresos, incapacidad laboral).
%      - Daño extrapatrimonial: daño moral (sufrimiento psíquico) y daño a la persona
%        (daño al proyecto de vida, daño biológico) — distinción Fernández Sessarego.
%   2. Analizar el daño iatrogénico (causado por el tratamiento) y su diferencia
%      con el riesgo propio de la enfermedad (daño inevitable vs. evitable).
%   3. Examinar el nexo causal como elemento crítico:
%      - Teoría de la causa adecuada (art. 1985 CC).
%      - Causalidad concurrente (culpa del médico + riesgo propio de la enfermedad).
%      - Pérdida de chance: ¿es aplicable en el sistema peruano?
%   4. Criterios de cuantificación del daño en la jurisprudencia peruana:
%      ausencia de baremo, discrecionalidad judicial, estándares emergentes.
%
% Referencias clave: Espinoza Espinoza, Fernández Sessarego, art. 1985 CC, CAS. relevantes
% Extensión estimada: 1.500–2.000 palabras

\section{El daño médico: tipología y criterios jurisprudenciales de valoración}
\label{sec:dano_medico}

